% Ad Familiares 13.6 --- headnote
% ad-familiares-13-06, latter part of 56 BC, written at Rome

\textit{Cicero to Quintus Valerius Orca, propraetor of Africa,
written at Rome in the latter part of 56 BC. Orca had crossed
to his province in the autumn; Cicero had seen him off
\textit{paludatum}, in the general's red cloak that marked the
formal departure for a command. The letter is the first of a
series Cicero is sending Orca on behalf of the Roman business
interests of his friend Publius Cuspius, twice principal of a
\textit{societas publicanorum} in Africa.}

\textit{The architecture of the letter is unusual: it is the
master commendation that legitimises the others. ``In this
letter I shall set out the grounds; in the rest I shall do
only what is needed to attach the agreed mark of recognition
and indicate that the man is a friend of Cuspius.'' The
particular man named here, Lucius Iulius, is being singled
out for a commendation pressed on Cicero with such zeal that
he is driven to a half-joke about the rhetorical art: Cuspius
expects something \textit{extraordinary}, and Cicero asks
Orca to make sure that, in the man's reception, the
commendation looks as if it had drawn upon the inmost
recesses of the rhetorical \textit{ars}. The closing aside on
the \textit{vultus} --- the welcome that the look of a
governor in his province can give --- is the working orator's
note to the working magistrate.}
